
\documentclass{article}
\usepackage{colortbl}
\usepackage{makecell}
\usepackage{multirow}
\usepackage{supertabular}

\begin{document}

\newcounter{utterance}

\centering \large Interaction Transcript for game `dond', experiment `semi', episode 3 with glm{-}5.2.
\vspace{24pt}

{ \footnotesize  \setcounter{utterance}{1}
\setlength{\tabcolsep}{0pt}
\begin{supertabular}{c@{$\;$}|p{.15\linewidth}@{}p{.15\linewidth}p{.15\linewidth}p{.15\linewidth}p{.15\linewidth}p{.15\linewidth}}
    \# & $\;$A & \multicolumn{4}{c}{Game Master} & $\;\:$B\\
    \hline

    \theutterance \stepcounter{utterance}  
    & & \multicolumn{4}{p{0.6\linewidth}}{
        \cellcolor[rgb]{0.9,0.9,0.9}{
            \makecell[{{p{\linewidth}}}]{
                \texttt{\tiny{[P1$\langle$GM]}}
                \texttt{Igrate igro sodelovalnega pogajanja, v kateri se morate dogovoriti s drugim igralcem o tem, kako razdeliti nabor predmetov med seboj.} \\
\\ 
\texttt{Pravila:} \\
\texttt{(a) Ti in drugi igralec dobite nabor predmetov. Vsakemu od vas je dana tudi tajna funkcija vrednosti, ki predstavlja, kako zelo cenite vsako vrsto predmeta.} \\
\texttt{(b) Izmenjujete si sporočila z drugim igralcem, da se dogovorite, kdo dobi katere predmete. Vsak lahko pošlje največ 5 sporočil ali zaključi igro prej tako, da v katerem koli trenutku daš skritega predloga.} \\
\texttt{(c) Prosimo vas, da pošljete skritega predloga, ki navaja predmete, ki jih želite, oblikovan v oglatih oglatih oklepajih, kot sledi: "{[}Predlog: <število> <ime predmeta>, <število> <ime predmeta>, <...>{]}"} \\
\texttt{(d) Če so vaši predlogi dopolnilni, tj. je dovolj predmetov za izpolnitev obeh predlogov, vsak igralec dobi rezultat na podlagi vsote vrednosti predmetov, ki jih je prejel. V nasprotnem primeru oba dobita nič točk.} \\
\texttt{(e) Vaš cilj je povečati vsoto vašega rezultata in rezultata drugega igralca.} \\
\\ 
\texttt{Začnimo.} \\
\\ 
\texttt{Niz razpoložljivih predmetov je:} \\
\\ 
\texttt{1 klešče, 2 opeke, 1 prstan, 1 ura} \\
\\ 
\texttt{Vaša tajna funkcija vrednosti je:} \\
\\ 
\texttt{klešče: 1, opeka: 0, prstan: 2, ura: 7} \\
\\ 
\texttt{POMEMBNO: Vaša sporočila, razen če gre za skriti predlog, se posredujejo neposredno drugemu igralcu, zato ne vključujte nobenega odgovora na pravila ali besedilo, ki napoveduje vaše sporočilo. Če želite podati skriti predlog, uporabite navedeni format. Ne uporabljajte oglatih oglatih oklepajev pri komunikaciji z drugim igralcem, saj bo to razumljeno kot vaš skritega predloga.} \\
            }
        }
    }
    & \\ \\

    \theutterance \stepcounter{utterance}  
    & \multicolumn{4}{p{0.6\linewidth}}{
        \cellcolor[rgb]{0.9,0.9,0.9}{
            \makecell[{{p{\linewidth}}}]{
                \texttt{\tiny{[P1$\rangle$GM]}}
                \texttt{Živjo! Predlagam, da začniva z razpravo o tem, kako vrednotiva predmete. Za začetek lahko povem, da me najbolj zanimata ura in prstan, medtem ko mi opeke ne pomenijo veliko. Kateri predmeti so najpomembnejši zate?} \\
            }
        }
    }
    & & \\ \\

    \theutterance \stepcounter{utterance}  
    & & & \multicolumn{2}{p{0.3\linewidth}}{
        \cellcolor[rgb]{0.9,0.9,0.9}{
            \makecell[{{p{\linewidth}}}]{
                \texttt{\tiny{[GM$|$GM]}}
                \texttt{continue} \\
            }
        }
    }
    & & \\ \\

    \theutterance \stepcounter{utterance}  
    & & & \multicolumn{2}{p{0.3\linewidth}}{
        \cellcolor[rgb]{0.9,0.9,0.9}{
            \makecell[{{p{\linewidth}}}]{
                \texttt{\tiny{[GM$|$GM]}}
                \texttt{Živjo! Predlagam, da začniva z razpravo o tem, kako vrednotiva predmete. Za začetek lahko povem, da me najbolj zanimata ura in prstan, medtem ko mi opeke ne pomenijo veliko. Kateri predmeti so najpomembnejši zate?} \\
            }
        }
    }
    & & \\ \\

    \theutterance \stepcounter{utterance}  
    & & \multicolumn{4}{p{0.6\linewidth}}{
        \cellcolor[rgb]{0.9,0.9,0.9}{
            \makecell[{{p{\linewidth}}}]{
                \texttt{\tiny{[GM$\rangle$P2]}}
                \texttt{Igrate igro sodelovalnega pogajanja, v kateri se morate dogovoriti s drugim igralcem o tem, kako razdeliti nabor predmetov med seboj.} \\
\\ 
\texttt{Pravila:} \\
\texttt{(a) Ti in drugi igralec dobite nabor predmetov. Vsakemu od vas je dana tudi tajna funkcija vrednosti, ki predstavlja, kako zelo cenite vsako vrsto predmeta.} \\
\texttt{(b) Izmenjujete si sporočila z drugim igralcem, da se dogovorite, kdo dobi katere predmete. Vsak lahko pošlje največ 5 sporočil ali zaključi igro prej tako, da v katerem koli trenutku daš skritega predloga.} \\
\texttt{(c) Prosimo vas, da pošljete skritega predloga, ki navaja predmete, ki jih želite, oblikovan v oglatih oglatih oklepajih, kot sledi: "{[}Predlog: <število> <ime predmeta>, <število> <ime predmeta>, <...>{]}"} \\
\texttt{(d) Če so vaši predlogi dopolnilni, tj. je dovolj predmetov za izpolnitev obeh predlogov, vsak igralec dobi rezultat na podlagi vsote vrednosti predmetov, ki jih je prejel. V nasprotnem primeru oba dobita nič točk.} \\
\texttt{(e) Vaš cilj je povečati vsoto vašega rezultata in rezultata drugega igralca.} \\
\\ 
\texttt{Začnimo.} \\
\\ 
\texttt{Niz razpoložljivih predmetov je:} \\
\\ 
\texttt{1 klešče, 2 opeke, 1 prstan, 1 ura} \\
\\ 
\texttt{Vaša tajna funkcija vrednosti je:} \\
\\ 
\texttt{klešče: 1, opeka: 1, prstan: 4, ura: 3} \\
\\ 
\texttt{POMEMBNO: Vaša sporočila, razen če gre za skriti predlog, se posredujejo neposredno drugemu igralcu, zato ne vključujte nobenega odgovora na pravila ali besedilo, ki napoveduje vaše sporočilo. Če želite podati skriti predlog, uporabite navedeni format. Ne uporabljajte oglatih oglatih oklepajev pri komunikaciji z drugim igralcem, saj bo to razumljeno kot vaš skritega predloga.} \\
\\ 
\texttt{Živjo! Predlagam, da začniva z razpravo o tem, kako vrednotiva predmete. Za začetek lahko povem, da me najbolj zanimata ura in prstan, medtem ko mi opeke ne pomenijo veliko. Kateri predmeti so najpomembnejši zate?} \\
            }
        }
    }
    & \\ \\

    \theutterance \stepcounter{utterance}  
    & & & \multicolumn{4}{p{0.6\linewidth}}{
        \cellcolor[rgb]{0.9,0.9,0.9}{
            \makecell[{{p{\linewidth}}}]{
                \texttt{\tiny{[GM$\langle$P2]}}
                \texttt{Živjo! Razumeviva se. Najbolj cenim opeke, klešče pa so mi manj pomembne. Predlagam, da vzameš prstan in uro, jaz pa opeke in klešče. Kako ti zveni?} \\
            }
        }
    }
     \\ \\

    \theutterance \stepcounter{utterance}  
    & & & \multicolumn{2}{p{0.3\linewidth}}{
        \cellcolor[rgb]{0.9,0.9,0.9}{
            \makecell[{{p{\linewidth}}}]{
                \texttt{\tiny{[GM$|$GM]}}
                \texttt{continue} \\
            }
        }
    }
    & & \\ \\

    \theutterance \stepcounter{utterance}  
    & & & \multicolumn{2}{p{0.3\linewidth}}{
        \cellcolor[rgb]{0.9,0.9,0.9}{
            \makecell[{{p{\linewidth}}}]{
                \texttt{\tiny{[GM$|$GM]}}
                \texttt{Živjo! Razumeviva se. Najbolj cenim opeke, klešče pa so mi manj pomembne. Predlagam, da vzameš prstan in uro, jaz pa opeke in klešče. Kako ti zveni?} \\
            }
        }
    }
    & & \\ \\

    \theutterance \stepcounter{utterance}  
    & & \multicolumn{4}{p{0.6\linewidth}}{
        \cellcolor[rgb]{0.9,0.9,0.9}{
            \makecell[{{p{\linewidth}}}]{
                \texttt{\tiny{[P1$\langle$GM]}}
                \texttt{Živjo! Razumeviva se. Najbolj cenim opeke, klešče pa so mi manj pomembne. Predlagam, da vzameš prstan in uro, jaz pa opeke in klešče. Kako ti zveni?} \\
            }
        }
    }
    & \\ \\

    \theutterance \stepcounter{utterance}  
    & \multicolumn{4}{p{0.6\linewidth}}{
        \cellcolor[rgb]{0.9,0.9,0.9}{
            \makecell[{{p{\linewidth}}}]{
                \texttt{\tiny{[P1$\rangle$GM]}}
                \texttt{To zveni odlično! Jaz bom vzel prstan in uro, ti pa opeke in klešče. Dogovorjeno!} \\
\\ 
\texttt{{[}Predlog: 1 prstan, 1 ura{]}} \\
            }
        }
    }
    & & \\ \\

    \theutterance \stepcounter{utterance}  
    & & & \multicolumn{2}{p{0.3\linewidth}}{
        \cellcolor[rgb]{0.9,0.9,0.9}{
            \makecell[{{p{\linewidth}}}]{
                \texttt{\tiny{[GM$|$GM]}}
                \texttt{proposal} \\
            }
        }
    }
    & & \\ \\

    \theutterance \stepcounter{utterance}  
    & & & \multicolumn{2}{p{0.3\linewidth}}{
        \cellcolor[rgb]{0.9,0.9,0.9}{
            \makecell[{{p{\linewidth}}}]{
                \texttt{\tiny{[GM$|$GM]}}
                [0, 0, 1, 1]
            }
        }
    }
    & & \\ \\

    \theutterance \stepcounter{utterance}  
    & & \multicolumn{4}{p{0.6\linewidth}}{
        \cellcolor[rgb]{0.9,0.9,0.9}{
            \makecell[{{p{\linewidth}}}]{
                \texttt{\tiny{[GM$\rangle$P2]}}
                \texttt{Drugi igralec je predložil skrivni predlog. Zdaj ste na vrsti, da predložite svoj predlog.} \\
            }
        }
    }
    & \\ \\

    \theutterance \stepcounter{utterance}  
    & & & \multicolumn{4}{p{0.6\linewidth}}{
        \cellcolor[rgb]{0.9,0.9,0.9}{
            \makecell[{{p{\linewidth}}}]{
                \texttt{\tiny{[GM$\langle$P2]}}
                \texttt{{[}Predlog: 1 klešče, 2 opeke{]}} \\
            }
        }
    }
     \\ \\

    \theutterance \stepcounter{utterance}  
    & & & \multicolumn{2}{p{0.3\linewidth}}{
        \cellcolor[rgb]{0.9,0.9,0.9}{
            \makecell[{{p{\linewidth}}}]{
                \texttt{\tiny{[GM$|$GM]}}
                \texttt{proposal} \\
            }
        }
    }
    & & \\ \\

    \theutterance \stepcounter{utterance}  
    & & & \multicolumn{2}{p{0.3\linewidth}}{
        \cellcolor[rgb]{0.9,0.9,0.9}{
            \makecell[{{p{\linewidth}}}]{
                \texttt{\tiny{[GM$|$GM]}}
                [1, 2, 0, 0]
            }
        }
    }
    & & \\ \\

    \theutterance \stepcounter{utterance}  
    & & & \multicolumn{2}{p{0.3\linewidth}}{
        \cellcolor[rgb]{0.9,0.9,0.9}{
            \makecell[{{p{\linewidth}}}]{
                \texttt{\tiny{[GM$|$GM]}}
                [[0, 0, 1, 1], [1, 2, 0, 0]]
            }
        }
    }
    & & \\ \\

\end{supertabular}
}

\end{document}
